\chapter{Аналитическая часть}
%<пишем про тему и все, связанное с ней>

% мат описание

Матрица -- прямоугольная таблица из чисел, содержащая $N$ строк и $M$ столбцов (см.~\ref{matr}):
\begin{align}
	\label{matr}
	\begin{pmatrix}
		a_{11} & ... & a_{1m} \\
		... & a_{ij} & ... \\
		a_{n1} & ... & a_{nm}
	\end{pmatrix},
\end{align}
где позиция каждого элемента матрицы задаётся индексами $i$ и $j$, отвечающими за положение в строке и столбце сответственно~\cite{lit1}. 

\section{Стандартный алгорим умножения матриц}

В стандартном (классическом) алгоритме умножения матриц каждый член результирующей матрицы является скалярным произведением элемента строки первой матрицы и соответствующего элемента столбца второй. Таким образом, умножение двух матриц размера $N$x$M$ и $M$x$Q$ имеет вид (см.~\ref{std_mul}):
\begin{align}
	\label{std_mul}
	\begin{pmatrix}
		a_{11} & ... & a_{1m} \\
		 & ... &  \\
		a_{n1} & ... & a_{nm}
	\end{pmatrix}
	\begin{pmatrix}
		b_{11} & ... & b_{1q} \\
		& ... &  \\
		b_{m1} & ... & b_{mq}
	\end{pmatrix}
	=
	\begin{pmatrix}
		c_{11} & ... & c_{1q} \\
		 & ... &  \\
		c_{n1} & ... & c_{nq}
	\end{pmatrix},
\end{align}
где каждый элемент результирующей матрицы имеет вид (см.~\ref{std_mul_el})~\cite{lit2}:
\begin{align}
	\label{std_mul_el}
	c_{ij} = \sum_{i=0}^{M} a_{ik}*b_{kj}.
\end{align}

\section{Алгорим Винограда}

Метод Винограда вносит ряд улучшений в стандартный алгоритм. Так как скалярное произведение двух векторов ${u(u1, u2, u3, u4)}$ и ${v(v1, v2, v3, v4)}$, имеющее вид (см.~\ref{std_vmul}):\\ 
\begin{align}
	\label{std_vmul}
	{u*v = u1*v1+u2*v2+u3*v3+u4*v4}
\end{align}
можно переписать как (см.~\ref{win_mul}):\\
\begin{align}
	\label{win_mul}
	{u*v = (u1+v2)*(u2+v1)+(u3+v4)*(u4+v3)-u1*u2-u3*u4-v1*v2-v3*v4}.
\end{align} 
Таким образом, часть значений можно вычислить заранее, что сокращает количество операций умножения~\cite{lit3}. 

\section{Алгорим Винограда с оптимизацией}
%#12539

Исходный алгоритм Винограда можно дополнительно оптимизировать.\\
Задание по варианту:
\begin{itemize}
	\item снять индекс в 1-ой части алгоритма;
	\item заменить $k++$ на $k \mathrel{+}= 2$.
\end{itemize}
Первое подзадание предполагает избавление от массива произведений пар элементов первой матрицы. Для этого можно вычитать эти произведения из соответствующих элементов результирующей матрицы в цикле с её заполнением, уменьшая тем количество циклов.\\
Замена $k++$ на $k \mathrel{+}= 2$ уменьшает количество умножений, так как теперь обращение к элементам выглядит как $matr[k][i]$, а не $matr[2k][i]$.

\section*{Вывод}
%<что сделали в аналите, кратко>

В данном разделе мной был проведён анализ теоретических материалов, объясняющих принцип работы алгоритмов умножения матриц, а также оптимизирован алгоритм Винограда согласно варианту.

\clearpage
